% Drop-in for Section VI-A (Scalability).
% Usage: % Drop-in for Section VI-A (Scalability).
% Usage: % Drop-in for Section VI-A (Scalability).
% Usage: % Drop-in for Section VI-A (Scalability).
% Usage: \input{scalability} inside \section{Applications and experiments},
%   e.g. before or after \input{deltasigma}.
% Before compiling in the main paper: copy
%   experiments/paper_experiments/11_stress_test/results/11_stress_ex123.pdf
% into the paper's fig/ directory, and add fig/ to \graphicspath in main.tex.

\rnew{
\subsection{Scalability}\label{subsec:scalability}
}

\begin{figure}[t]
  \centering
  \includegraphics[width=\columnwidth]{11_stress_ex123.pdf}
  \caption{\rnew{Stress-test timing for three TREs ($K=20$ samples per
      data point, mean sampling time in seconds).
      Bars decompose into one-time volume-computation cost and per-sample cost.
      \emph{Top:} $\eex{1}=(\trest{a}{\leq 1})^*$.
      \emph{Middle:} $\eex{2}=a^*\concat\trest{a^*}{[1,2]}$ ($T=2.0$).
      \emph{Bottom:} $\eex{3}=(\trest{aa}{\leq 1}\concat a)
        \cap(a\concat\trest{aa}{\leq 1})$ ($n=3$).}}
  \label{fig:stress}
\end{figure}

\rnew{%
Fig.~\ref{fig:stress} characterises scaling across three TREs, each
stressing a different aspect of the algorithms ($K=20$ samples per data point).
Timing decomposes into a \emph{one-time volume computation}
(Thm.~\ref{thm:polynomials}), paid once per specification regardless
of how many words are subsequently drawn, and a \emph{per-sample} cost.
In typical use, the former is amortised across all samples.

For the unambiguous expression $\eex{1}$, Alg.~\ref{alg:slicesampling}
proceeds without rejection and scales to $n=22$ with cost dominated
by volume computation.
For the ambiguous $\eex{2}$, volume computation remains inexpensive, while
per-sample cost grows with~$n$, consistent with Alg.~\ref{alg:smartsampling}
incurring an increasing density correction cost as the average multiplicity of
generated words grows with~$n$ (Thm.~\ref{thm:slicesampling}).
For the intersection expression $\eex{3}$, Alg.~\ref{alg:gensmartsampling}
draws candidates from the first subexpression and accepts those also belonging
to the second. As $T$ approaches the feasibility boundary at $T=2$ from
below, the ratio of the intersection volume to the enclosing language volume
vanishes, so the expected cost per accepted sample grows accordingly.
This is the expected behaviour of rejection sampling as the
acceptance probability decreases.

The density correction cost in Alg.~\ref{alg:smartsampling} can be significant for
highly ambiguous expressions. The expected number of trials is characterised
by the multiplicity structure of the expression (Thm.~\ref{thm:slicesampling}).
The rejection-based handling of intersection in Alg.~\ref{alg:gensmartsampling}
is a known limitation of the approach: one can construct TREs with arbitrarily
thin intersections, and the cost scales with the inverse of the acceptance
probability.
}
 inside \section{Applications and experiments},
%   e.g. before or after \input{deltasigma}.
% Before compiling in the main paper: copy
%   experiments/paper_experiments/11_stress_test/results/11_stress_ex123.pdf
% into the paper's fig/ directory, and add fig/ to \graphicspath in main.tex.

\rnew{
\subsection{Scalability}\label{subsec:scalability}
}

\begin{figure}[t]
  \centering
  \includegraphics[width=\columnwidth]{11_stress_ex123.pdf}
  \caption{\rnew{Stress-test timing for three TREs ($K=20$ samples per
      data point, mean sampling time in seconds).
      Bars decompose into one-time volume-computation cost and per-sample cost.
      \emph{Top:} $\eex{1}=(\trest{a}{\leq 1})^*$.
      \emph{Middle:} $\eex{2}=a^*\concat\trest{a^*}{[1,2]}$ ($T=2.0$).
      \emph{Bottom:} $\eex{3}=(\trest{aa}{\leq 1}\concat a)
        \cap(a\concat\trest{aa}{\leq 1})$ ($n=3$).}}
  \label{fig:stress}
\end{figure}

\rnew{%
Fig.~\ref{fig:stress} characterises scaling across three TREs, each
stressing a different aspect of the algorithms ($K=20$ samples per data point).
Timing decomposes into a \emph{one-time volume computation}
(Thm.~\ref{thm:polynomials}), paid once per specification regardless
of how many words are subsequently drawn, and a \emph{per-sample} cost.
In typical use, the former is amortised across all samples.

For the unambiguous expression $\eex{1}$, Alg.~\ref{alg:slicesampling}
proceeds without rejection and scales to $n=22$ with cost dominated
by volume computation.
For the ambiguous $\eex{2}$, volume computation remains inexpensive, while
per-sample cost grows with~$n$, consistent with Alg.~\ref{alg:smartsampling}
incurring an increasing density correction cost as the average multiplicity of
generated words grows with~$n$ (Thm.~\ref{thm:slicesampling}).
For the intersection expression $\eex{3}$, Alg.~\ref{alg:gensmartsampling}
draws candidates from the first subexpression and accepts those also belonging
to the second. As $T$ approaches the feasibility boundary at $T=2$ from
below, the ratio of the intersection volume to the enclosing language volume
vanishes, so the expected cost per accepted sample grows accordingly.
This is the expected behaviour of rejection sampling as the
acceptance probability decreases.

The density correction cost in Alg.~\ref{alg:smartsampling} can be significant for
highly ambiguous expressions. The expected number of trials is characterised
by the multiplicity structure of the expression (Thm.~\ref{thm:slicesampling}).
The rejection-based handling of intersection in Alg.~\ref{alg:gensmartsampling}
is a known limitation of the approach: one can construct TREs with arbitrarily
thin intersections, and the cost scales with the inverse of the acceptance
probability.
}
 inside \section{Applications and experiments},
%   e.g. before or after \input{deltasigma}.
% Before compiling in the main paper: copy
%   experiments/paper_experiments/11_stress_test/results/11_stress_ex123.pdf
% into the paper's fig/ directory, and add fig/ to \graphicspath in main.tex.

\rnew{
\subsection{Scalability}\label{subsec:scalability}
}

\begin{figure}[t]
  \centering
  \includegraphics[width=\columnwidth]{11_stress_ex123.pdf}
  \caption{\rnew{Stress-test timing for three TREs ($K=20$ samples per
      data point, mean sampling time in seconds).
      Bars decompose into one-time volume-computation cost and per-sample cost.
      \emph{Top:} $\eex{1}=(\trest{a}{\leq 1})^*$.
      \emph{Middle:} $\eex{2}=a^*\concat\trest{a^*}{[1,2]}$ ($T=2.0$).
      \emph{Bottom:} $\eex{3}=(\trest{aa}{\leq 1}\concat a)
        \cap(a\concat\trest{aa}{\leq 1})$ ($n=3$).}}
  \label{fig:stress}
\end{figure}

\rnew{%
Fig.~\ref{fig:stress} characterises scaling across three TREs, each
stressing a different aspect of the algorithms ($K=20$ samples per data point).
Timing decomposes into a \emph{one-time volume computation}
(Thm.~\ref{thm:polynomials}), paid once per specification regardless
of how many words are subsequently drawn, and a \emph{per-sample} cost.
In typical use, the former is amortised across all samples.

For the unambiguous expression $\eex{1}$, Alg.~\ref{alg:slicesampling}
proceeds without rejection and scales to $n=22$ with cost dominated
by volume computation.
For the ambiguous $\eex{2}$, volume computation remains inexpensive, while
per-sample cost grows with~$n$, consistent with Alg.~\ref{alg:smartsampling}
incurring an increasing density correction cost as the average multiplicity of
generated words grows with~$n$ (Thm.~\ref{thm:slicesampling}).
For the intersection expression $\eex{3}$, Alg.~\ref{alg:gensmartsampling}
draws candidates from the first subexpression and accepts those also belonging
to the second. As $T$ approaches the feasibility boundary at $T=2$ from
below, the ratio of the intersection volume to the enclosing language volume
vanishes, so the expected cost per accepted sample grows accordingly.
This is the expected behaviour of rejection sampling as the
acceptance probability decreases.

The density correction cost in Alg.~\ref{alg:smartsampling} can be significant for
highly ambiguous expressions. The expected number of trials is characterised
by the multiplicity structure of the expression (Thm.~\ref{thm:slicesampling}).
The rejection-based handling of intersection in Alg.~\ref{alg:gensmartsampling}
is a known limitation of the approach: one can construct TREs with arbitrarily
thin intersections, and the cost scales with the inverse of the acceptance
probability.
}
 inside \section{Applications and experiments},
%   e.g. before or after \input{deltasigma}.
% Before compiling in the main paper: copy
%   experiments/paper_experiments/11_stress_test/results/11_stress_ex123.pdf
% into the paper's fig/ directory, and add fig/ to \graphicspath in main.tex.

\rnew{
\subsection{Scalability}\label{subsec:scalability}
}

\begin{figure}[t]
  \centering
  \includegraphics[width=\columnwidth]{11_stress_ex123.pdf}
  \caption{\rnew{Stress-test timing for three TREs ($K=20$ samples per
      data point, mean sampling time in seconds).
      Bars decompose into one-time volume-computation cost and per-sample cost.
      \emph{Top:} $\eex{1}=(\trest{a}{\leq 1})^*$.
      \emph{Middle:} $\eex{2}=a^*\concat\trest{a^*}{[1,2]}$ ($T=2.0$).
      \emph{Bottom:} $\eex{3}=(\trest{aa}{\leq 1}\concat a)
        \cap(a\concat\trest{aa}{\leq 1})$ ($n=3$).}}
  \label{fig:stress}
\end{figure}

\rnew{%
Fig.~\ref{fig:stress} characterises scaling across three TREs, each
stressing a different aspect of the algorithms ($K=20$ samples per data point).
Timing decomposes into a \emph{one-time volume computation}
(Thm.~\ref{thm:polynomials}), paid once per specification regardless
of how many words are subsequently drawn, and a \emph{per-sample} cost.
In typical use, the former is amortised across all samples.

For the unambiguous expression $\eex{1}$, Alg.~\ref{alg:slicesampling}
proceeds without rejection and scales to $n=22$ with cost dominated
by volume computation.
For the ambiguous $\eex{2}$, volume computation remains inexpensive, while
per-sample cost grows with~$n$, consistent with Alg.~\ref{alg:smartsampling}
incurring an increasing density correction cost as the average multiplicity of
generated words grows with~$n$ (Thm.~\ref{thm:slicesampling}).
For the intersection expression $\eex{3}$, Alg.~\ref{alg:gensmartsampling}
draws candidates from the first subexpression and accepts those also belonging
to the second. As $T$ approaches the feasibility boundary at $T=2$ from
below, the ratio of the intersection volume to the enclosing language volume
vanishes, so the expected cost per accepted sample grows accordingly.
This is the expected behaviour of rejection sampling as the
acceptance probability decreases.

The density correction cost in Alg.~\ref{alg:smartsampling} can be significant for
highly ambiguous expressions. The expected number of trials is characterised
by the multiplicity structure of the expression (Thm.~\ref{thm:slicesampling}).
The rejection-based handling of intersection in Alg.~\ref{alg:gensmartsampling}
is a known limitation of the approach: one can construct TREs with arbitrarily
thin intersections, and the cost scales with the inverse of the acceptance
probability.
}
